\section{Filipino Orthography and Phonology}

Due to the introduction of loanwords from Spanish, English, and other languages,
the phonological and orthographical systems of the Filipno language have evolved in complexity to 
adapt for growing phonemes and graphemes. Originally, Baybayin script was used as
syllabary for the writing of Filipino. Post-introduction of the Latin alphabet during 
the Spanish colonial era, local orthographic and phonetic inventories expanded to
emphasize contrasts between newly introduced vowel and consonant sounds. With Baybayin
having shared symbols for \textit{(/e/)} and \textit{(/i/)}, and \textit{(/o/)} and 
\textit{(/u/)}, Filipino utilized a three-vowel system of \textit{(/a, i, u/)}, with the
shared symbols being allophones~\cite{schachter1972}. But to distinguish between Spanish
loanwords, the three vowel symbols were split into five distinct letters 
\textit{(/a, e, i, o u/)}~\cite{almario2014}. Native consonant phoneme and grapheme
structures also evolved to further accomodate additional sounds introduced through
borrowing, such as \textit{/f, v, z, \textipa{S}, \textipa{dZ}, \textipa{tS}/}. 
The introduction of new phonetic and orthographic variants create adapted variations 
of these phonemes to match local phonetics~\cite{llamzon1966}, for example \textit{/\textipa{tS}/} 
and \textit{/\textipa{dZ}/}, are adapted into the language as \textit{/ts/} and \textit{/dy/}
~\cite{llamzon1966}.

Considering that Filipino had roots in syllabary, it is a syllable-timed language, 
where syllables are the fundamental structural unit and underlying framework for the 
distribution of phonemes in the language~\cite{llamzon1966, schachter1972}. 
Filipino has a basic syllable structure of consonant-vowel-consonant (CVC), but has 
evolved in complexity due to introductions of loanwords (\textit{ex. komiks (CVCC), trump, (CCVCC)}).
While adaptation of loanwords into Filipino allow it to mostly retain its original syllable 
structure, speakers can change the phonological structure of a loanword (vowel insertion,
sound deletion) to better match the native phonetic system.

In modern Filipino, the phonological features of English loanwords are reflected in 
adapted speech and text, and these are often subject to varying degrees of adaptation
depending on a speaker's education, socioeconomic status, and exposure to English media.
This sociolinguistic phenomenon can be categorized by the sociolectal identities of acrolect, 
mesolect, and basilect~\cite{tayao2004} in terms of their linguistic proximity to General American English 
(gAmE).